\documentclass[aspectratio=169,11pt]{beamer}

\usepackage[T1]{fontenc}
\usepackage[utf8]{inputenc}
\usepackage{lmodern}
\usepackage{amsmath,amssymb}
\usepackage{booktabs}
\usepackage{xcolor}
\usepackage{hyperref}
\usepackage{microtype}

\usetheme{Madrid}
\usecolortheme{default}

\definecolor{LaurierPurple}{HTML}{4B116F}
\definecolor{DeepBlue}{HTML}{0047AB}

\setbeamertemplate{navigation symbols}{}
\setbeamercolor{title}{fg=white,bg=LaurierPurple}
\setbeamercolor{frametitle}{fg=white,bg=LaurierPurple}
\setbeamercolor{block title}{fg=white,bg=DeepBlue}
\setbeamercolor{structure}{fg=DeepBlue}

\hypersetup{colorlinks=true,urlcolor=DeepBlue,linkcolor=DeepBlue}

\title[Trees \& SHAP]{Module 5: Tree-Based Methods, Classification, and Interpretability}
\subtitle{EC 410K / EC 610I --- Machine Learning in Economics}
\author{M. Jahangir Alam, PhD}
\institute{Wilfrid Laurier University}
\date{Spring 2026}

\begin{document}

\begin{frame}[plain]
  \titlepage
\end{frame}

\begin{frame}{Module overview}
  \begin{itemize}
    \item Tree ensembles (random forests, boosting) are workhorses for prediction and classification in applied work.
    \item They capture nonlinearities and interactions but can be harder to interpret than linear models.
    \item SHAP provides a principled way to attribute predictions to features---useful, but not a substitute for causal reasoning.
  \end{itemize}
\end{frame}

\begin{frame}{Learning objectives}
  \begin{enumerate}
    \item Explain bagging (random forests) vs boosting at intuition level.
    \item Train/evaluate a classification model with out-of-sample AUC.
    \item Interpret SHAP summary plots for global and local insight.
    \item Connect to Lundberg and Lee (2017) and Lundberg et al.\ (2020).
  \end{enumerate}
\end{frame}

\begin{frame}{Trees: idea}
  \begin{itemize}
    \item Recursive partitioning: split feature space to minimize impurity (Gini, entropy) or loss.
    \item Pros: flexible; interactions; minimal tuning sometimes.
    \item Cons: can overfit; unstable individually; extrapolation weaknesses.
  \end{itemize}
\end{frame}

\begin{frame}{Random forests (bagging + random feature subsets)}
  \begin{itemize}
    \item Average many deep trees trained on bootstrap samples.
    \item Variance reduction; improved generalization vs a single tree.
  \end{itemize}
\end{frame}

\begin{frame}{Boosting (high level)}
  \begin{itemize}
    \item Sequentially fit weak learners to correct residuals.
    \item Often strong accuracy; requires more careful tuning and can overfit.
  \end{itemize}
\end{frame}

\begin{frame}{Classification metrics economists use}
  \begin{itemize}
    \item Accuracy can mislead with imbalance; ROC-AUC summarizes discrimination.
    \item Calibration matters for policy: predicted probabilities should match frequencies when possible.
  \end{itemize}
\end{frame}

\begin{frame}{SHAP values: cooperative game interpretation}
  \begin{itemize}
    \item Shapley values from game theory allocate ``payout'' (prediction) across features.
    \item Local explanation: why did this unit get \(\widehat{p}\)?
    \item Global view: aggregate \(|\text{SHAP}|\) for variable importance comparisons.
  \end{itemize}
\end{frame}

\begin{frame}{SHAP \(\neq\) causality}
  \begin{alertblock}{Common mistake}
    ``SHAP says \(X_1\) matters most, therefore \(X_1\) is a policy lever.''
  \end{alertblock}
  \begin{itemize}
    \item SHAP explains the fitted model; confounding and endogeneity remain.
  \end{itemize}
\end{frame}

\begin{frame}{Replication paper connections}
  \begin{itemize}
    \item \textbf{Goldin et al.\ (2016):} classify plan types or predict enrollment responses; interpret drivers via SHAP (prediction narrative).
    \item \textbf{Hansen (2015):} predict recidivism vs estimate deterrence at threshold---keep tasks separate.
  \end{itemize}
\end{frame}

\begin{frame}{Python / Colab demonstration}
  \begin{enumerate}
    \item Simulate a logistic data-generating process with interactions/nonlinearities.
    \item Train random forest; report holdout AUC.
    \item Compute SHAP with TreeExplainer; show summary plot.
  \end{enumerate}
\end{frame}

\begin{frame}[fragile]{Colab setup}
  \begin{verbatim}
!pip -q install numpy pandas scikit-learn matplotlib shap
  \end{verbatim}
\end{frame}

\begin{frame}{Student / group assignment}
  \begin{itemize}
    \item Explain one example where feature importance differs from SHAP attributions and why.
    \item Present Lundberg and Lee (2017) or tree SHAP paper (2020).
    \item Run Colab; discuss computational cost and subsampling strategies.
    \item Propose an interpretability study for a replication paper that stays within defensible claims.
  \end{itemize}
\end{frame}

\begin{frame}{Summary}
  \begin{itemize}
    \item Tree ensembles are strong predictors; interpretability tools help audit and communicate.
    \item SHAP is useful for explaining model behavior, not establishing causal mechanisms.
    \item Always pair interpretability with identification for policy claims.
  \end{itemize}
\end{frame}

\begin{frame}[plain]
  \begin{center}
    \Large Questions?\\[0.8em]
    \normalsize Next: forecasting with ML/deep learning (Module 6).
  \end{center}
\end{frame}

\end{document}
